\chapter{Integrali di linea di $2^\circ$ specie}

\section{Definizione di Campo Vettoriale}
Un \textbf{campo vettoriale} è una funzione:
$$\vec{F}: D \subset \mathbb{R}^n \rightarrow \mathbb{R}^n$$
che associa a ogni punto del dominio $D$ un vettore di $\mathbb{R}^n$.

Si può scrivere nella forma:
$$\vec{F} = F_1 \vec{e}_1 + \dots + F_n \vec{e}_n$$
dove $\hat{e}_i$ sono i versori della base canonica e $F_i$ sono le componenti del campo vettoriale.

\section{Definizione di Integrale di linea di $2^\circ$ specie}
Dati:
\begin{itemize}
    \item Un campo vettoriale $\vec{F}: D \subset \mathbb{R}^n \rightarrow \mathbb{R}^n$.
    \item Una \textbf{curva orientata} $\gamma$ (cioè una curva su cui abbiamo fissato un verso di percorrenza).
    \item Il versore tangente $\vec{t}$ (con $\|\vec{t}\| = 1$), avente verso concorde all'orientamento scelto per la curva.
\end{itemize}

L'integrale di linea di $2^\circ$ specie del campo vettoriale $\vec{F}$ lungo la curva $\gamma$ è definito come l'integrale di prima specie della sua componente tangenziale:
$$\int_{\gamma} \vec{F} \cdot \vec{t} \, ds$$

\section{Come si calcola}
Sia $\vec{r}(t)$ con $t \in [a,b]$ una parametrizzazione regolare di $\gamma$, concorde con l'orientamento della curva. 
Sappiamo che il versore tangente è dato da:
$$\vec{t} = \frac{\frac{d\vec{r}}{dt}}{\left\| \frac{d\vec{r}}{dt} \right\|}$$
e l'elemento d'arco è:
$$dl = \left\| \frac{d\vec{r}}{dt} \right\| dt$$

Sostituendo nell'integrale otteniamo:
$$\int_{\gamma} \vec{F} \cdot \vec{t} \, ds = \int_{a}^{b} \vec{F}(t) \cdot \frac{\frac{d\vec{r}}{dt}}{\left\| \frac{d\vec{r}}{dt} \right\|} \left\| \frac{d\vec{r}}{dt} \right\| dt = \int_{a}^{b} \vec{F}(t) \cdot \frac{d\vec{r}}{dt} \, dt$$

Esplicitando il prodotto scalare in termini di componenti, l'integrale si calcola come:
$$\int_{a}^{b} \left( F_1(x_1(t), \dots, x_n(t))\frac{dx_1}{dt} + \dots + F_n(x_1(t), \dots, x_n(t))\frac{dx_n}{dt} \right) dt$$

Questa espressione viene spesso scritta in modo compatto utilizzando la notazione delle \textbf{forme differenziali}:
$$\int_{\gamma} F_1 \, dx_1 + \dots + F_n \, dx_n \, dt$$

